<# Copyright (c) 2026 Benjamin Benno Falkner, SPDX-License-Identifier: MIT #>
<# Q-Phase Exam #> 
\documentclass[lang=ngerman]{edoxarticle}

<% if node.args.macros | default(value=[]) | length > 0 %>

<% for m in node.args.macros %>
\<% if m.provide %>providecommand<% else %>newcommand<% endif %>{\<< m.name >>}<% if m.args != null %>[<< m.args >>]<% endif %><% if m.default != null %>[<< m.default >>]<% endif %>{<< m.body >>}
<% endfor %>
<% endif %>
<% if node.args.examName != null %>
\renewcommand\examname{<< node.args.examName >>}
<% endif %>
<% if node.args.pointsName != null %>
\renewcommand\pointsname{<< node.args.pointsName >>}
<% endif %>

\begin{document}
<% if node.args.date != null %>
\pdate{<< node.args.date >>}
<% endif %>
\begin{exam}<% if node.args.title != null or node.args.titleTex != null %>{<< components.helpers.title(node=node) >>}<% endif %>

\vskip20pt
{\LARGE\textbf{Vorblatt}}
\vskip20pt

Für die Bearbeitung der Aufgaben des Prüfungsteils A sind Zeichengeräte und ein
Wörterbuch zur deutschen Rechtschreibung zugelassen. Eine Verwendung von
weiteren Hilfsmitteln ist nicht zulässig.

Sie müssen den Prüfungsteil A \textbf{spätestens << node.args.timePartA | default(value=35) >> Minuten} nach Prüfungsbeginn
abgeben. Die Hilfsmittel Taschenrechner und Formelsammlung erhalten Sie nach
Abgabe von Prüfungsteil A.

Sie müssen die \textbf{drei Pflichtaufgaben und zwei Wahlpflichtaufgaben} im
Prüfungsteil A bearbeiten.

\textbf{Tragen Sie die Aufgabennummern der zwei Wahlpflichtaufgaben ein, die
bewertet werden sollen.}

\vskip20pt
Wahlpflichtaufgabe Nr. \hbox to 4cm{\hrulefill}
\vskip10pt
Wahlpflichtaufgabe Nr. \hbox to 4cm{\hrulefill}
\vskip30pt
\textbf{Bestätigen Sie Ihre Entscheidung mit Ihrer Unterschrift.}
\vskip10pt
\hbox to \textwidth{Name, Vorname des Prüflings: \kern3.4cm\leaders\hrule\hfill}\vskip10pt
\hbox to \textwidth{Unterschrift des Prüflings:  \kern4cm\leaders\hrule\hfill}

\newpage

<< components.helpers.blocks(nodes=node.children) >>
<# \setpoints is needed when only some exercises are graded #>

\vskip20pt
\textbf{Zugelassene Hilfsmittel für Prüfungsteil B:}
\begin{itemize}
\item CAS-Taschenrechner 
\item Formelsammlung
\item Wörterbuch zur deutschen Rechtschreibung
\end{itemize}


<% if node.args.totalPoints != null %>
\setpoints{<< node.args.totalPoints >>}
<% endif %>
\end{exam}
\end{document}